\section{Experimental Setup}

This section should make the evaluation protocol precise and reproducible.

\subsection{Datasets and Tasks}

Describe the dataset source, version, schema, split, annotation process, and task definition. If the paper introduces a new benchmark task, define the input, output, quality criteria, and expected use.

\subsection{Baselines}

Include strong and relevant baselines. Typical IRE baselines include:

\begin{itemize}
  \item A single-prompt LLM baseline.
  \item A retrieval-augmented single-agent baseline.
  \item A rule-based or traditional NLP baseline where applicable.
  \item A human-authored or expert-reviewed reference where available.
  \item ChatREQ ablations that remove knowledge, agent critique, graph context, or human-in-the-loop control.
\end{itemize}

\subsection{Metrics}

Define task-specific metrics such as completeness, consistency, conflict detection precision and recall, traceability quality, feasibility accuracy, human effort, and downstream utility.

\subsection{Ablations}

Plan ablations that test the contribution of knowledge sources, agent roles, retrieval, graph reasoning, human review policies, and workflow stages.

\subsection{Reproducibility}

Report the code version, dataset version, prompt version, knowledge base version, model configuration, random seeds where applicable, and scripts used to generate results.

